\begin{textsample}{POS Dim 1 – 2017\_c4 – Score -3.00 – 2017\_c4/enade\_2017\_eng\_flo\_617fb2f4dc.txt}  \label{ex:f1_pos_274}
CAPÍTULO 4 PERCEPÇÃO DA PROVA As análises feitas neste capítulo tratam das percepções dos concluintes da Área de Engenharia Florestal sobre a prova aplicada no Enade/2017.
Essas percepções foram mensuradas por meio de nove questões que avaliaram desde o grau de dificuldade da prova até o tempo gasto para concluí -la.
As percepções sobre a prova foram relacionadas ao desempenho dos estudantes e à Grande Região de funcionamento do curso.
O questionário de percepção da prova encontra -se ao final do Anexo VII, que traz a reprodução do exame.
O desempenho dos estudantes foi classificado em quatro quartos.
Para tanto, esse desempenho foi ordenado de forma ascendente.
O percentil 25, P25, também conhecido como primeiro quartil, é a nota de desempenho que deixa um quarto ( 25 \% ) dos valores observados abaixo e três quartos acima.
A Figura 1 apresenta uma ilustração desse conceito.
O quarto inferior de desempenho é composto pelas notas abaixo do primeiro quartil.
Já o percentil 75, P75, também conhecido como terceiro quartil, é o valor para o qual há três quartos ( 75 \% ) dos dados abaixo e um quarto acima dele.
O quarto superior de desempenho é composto pelas notas iguais ou acima do terceiro quartil.
O percentil 50, P50, também conhecido como mediana, é o valor que divide as notas em dois conjuntos de igual tamanho.
O segundo quarto inclui valores entre o primeiro quartil ( P25 ) e a mediana.
O terceiro quarto contém os valores entre a mediana ( P50 ) e o terceiro quartil ( P75 ).
Vale ressaltar que percentis, quartis e medianas são pontos que não obrigatoriamente pertencem ao conjunto original de dados, ao passo que os quartos são subconjuntos dos dados originais.
A seguir, serão apresentados gráficos com resultados selecionados, relativos às nove questões avaliadas por grupos de estudantes.
Os gráficos apresentam nas barras o percentual de alunos que assinalaram uma das opções ou a soma das porcentagens daqueles que assinalaram duas ( ou três ) delas.
Por exemplo, para as questões 1 e 2, os gráficos apresentam a porcentagem total de participantes que assinalaram as opções (D)Difícil e ( E ) Muito difícil.
Em cada barra foram assinalados também os extremos do intervalo de confiança de 95 \% como linhas verticais unidas por uma linha horizontal na forma da letra H maiúscula.
As Tabelas no Anexo II apresentam os valores absolutos e a distribuição percentual17 das alternativas válidas das nove questões, segundo o mesmo recorte de desempenho dos alunos e Grande Região de funcionamento do curso. 4.1 GRAU DE DIFICULDADE DA PROVA 4.1.1 Componente de Formação Geral Ao avaliarem Qual o grau de dificuldade desta prova na parte de Formação Geral? ( Questão 1 ), 19,8 \% do grupo de inscritos e presentes optaram pelas alternativas Difícil ou Muito difícil.
Entretanto, para mais da metade dos estudantes ( 66,2 \% ), o Componente de Formação Geral da prova foi considerado com grau de dificuldade Médio ( Gráfico 4.1, Gráfico 4.2 e, no Anexo II, a Tabela II.1 ).
O percentual de estudantes que consideraram a prova como Difícil ou Muito difícil foi maior na região Sul, onde a proporção foi de 24,9 \%, enquanto a menor proporção ocorreu na região Sudeste, com 16,0 \%.
No Gráfico 4.1, é possível observar que as diferenças entre a região Sul e as regiões Sudeste e Norte são estatisticamente significativas.
Nas Grandes Regiões, a proporção de presentes à prova que consideraram o Componente de Formação Geral como sendo de grau de dificuldade Médio esteve entre 59,9 \% na região Nordeste, e 71,8 \% na região Norte. 17 Cumpre lembrar uma das convenções para tabelas numéricas ( pág. iii ) sobre a possibilidade de a soma das partes não resultar em 100 \% por questões de arredondamento.
O percentual de alunos que consideraram a prova Difícil ou Muito difícil foi decrescente em função dos quartos de desempenho : 25,8 \% no primeiro quarto e 15,7 \% no quarto, grupo de melhor desempenho na prova.
Nos quartos de desempenho intermediários, a proporção de alunos que consideraram a prova Difícil ou Muito difícil foi de 19,0 \% no segundo quarto e de 18,6 \% no terceiro quarto.
As diferenças entre os quartos inferior e superior de desempenho são estatisticamente significativas.
Para todos os quartos de desempenho, a alternativa modal para essa pergunta foi Médio, com 62,6 \% e 67,2 \% dos respondentes nos quartos extremos, primeiro e quarto, respectivamente. 4.1.2 Componente de Conhecimento Específico Ao responderem à Questão 2 – Qual o grau de dificuldade desta prova na parte de Componente Específico? – 29,3 \% do grupo de estudantes classificaram -na como Difícil ou Muito difícil.
Além disso, o Componente de Conhecimento Específico da prova foi considerado com grau de dificuldade Médio por 62,4 \% dos alunos ( Gráfico 4.3, Gráfico 4.4, e, no Anexo II, a Tabela II.2 ).
A análise das respostas dos estudantes quanto ao grau de dificuldade do Componente de Conhecimento Específico da prova, agregado por Grande Região, mostra que a diferença entre a menor e as três maiores proporções de alunos que a avaliaram como Difícil ou Muito difícil é estatisticamente significativa : a menor na região Sudeste ( 22,5 \% ) e as três maiores proporções nas regiões Norte ( 31,0 \% ), Centro-Oeste ( 32,9 \% ) e Sul ( 33,2 \% ), e o percentual de alunos que classificaram o grau de dificuldade como Médio, no Componente de Conhecimento Específico, variou de 57,4 \% a 67,7 \%, nas regiões Sul e Sudeste, respectivamente.
Considerando -se a avaliação da dificuldade das questões do Componente de Conhecimento Específico da prova, de acordo com o desempenho dos estudantes, observa -se que há diferença estatisticamente significativa dos resultados entre o primeiro e os dois quartos de desempenho superiores, bem como entre o quarto superior e os dois quartos inferiores de desempenho.
A proporção dos que classificaram a parte específica como Difícil ou Muito difícil variou de 23,7 \% ( último quarto ) a 34,1 \% ( primeiro quarto ).
A alternativa modal para a Questão 2 foi o grau Médio, com 58,0 \% do quarto inferior e 65,8 \% do superior optando por essa resposta. 4.2 EXTENSÃO DA PROVA EM RELAÇÃO AO TEMPO TOTAL Indagados quanto à extensão da prova, em relação ao tempo total oferecido para a sua resolução ( Questão 3 ), os estudantes apontaram, com maior incidência, a alternativa que considerava a extensão adequada, para todas as agregações consideradas ( Gráfico 4.5, Gráfico 4.6, e, no Anexo II, a Tabela II.3 ).
O percentual de alunos que responderam ser a extensão da prova adequada foi de 68,8 \%.
Já 25,9 \% dos inscritos presentes consideraram que a prova foi longa ou muito longa, e 5,2 \% a avaliaram como curta ou muito curta.
Dentre as Grandes Regiões, a proporção daqueles que avaliaram a prova como longa ou muito longa, em relação ao tempo total destinado à sua resolução variou de 22,5 \% na região Norte até 29,5 \% na região Nordeste.
Não há diferença estatisticamente significativa dentre as regiões.
Considerando -se o desempenho dos alunos, nota -se ainda que 64,6 \% consideraram a extensão da prova adequada no quarto de desempenho inferior, e 71,2 \%, no de melhor desempenho ( quarto superior ).
Nos quartos intermediários, essa proporção foi de 68,3 \% no segundo quarto e de 71,2 \% no terceiro.
No Gráfico 4.6, verifica -se que dentre os estudantes que consideraram a prova longa ou muito longa, em função dos quartos de desempenho, as proporções foram de 29,6 \% no quarto inferior e de 24,1 \% no superior.
Observa -se, porém, que não há diferença estatisticamente significativa entre as proporções de escolha dessa resposta. 4.3 COMPREENSÃO DOS ENUNCIADOS DAS QUESTÕES 4.3.1 Componente de Formação Geral Em relação aos enunciados das questões do Componente de Formação Geral ( Questão 4 ), as opiniões foram positivas, já que 80,9 \% dos alunos avaliados consideraram todos ou a maioria dos enunciados das questões claros e objetivos ( Gráfico 4.7, Gráfico 4.8, e, no Anexo II, a Tabela II.4 ).
Na análise regional, a percentagem de estudantes que avaliaram que todos ou a maioria dos enunciados das questões do Componente de Formação Geral estavam claros e objetivos variou de 80,0 \% na região Sudeste a 81,8 \% na região Norte, não havendo diferenças estatisticamente significativas entre as regiões.
A análise das percepções dos estudantes sobre a clareza e objetividade dos enunciados permite afirmar que todos, ou a maioria dos enunciados de questões relativas ao Componente de Formação Geral, foram considerados claros e objetivos para a maior parte dos respondentes ( maior ou igual a 80,0 \% em todas as regiões e maior ou igual a 74,1 \% para todos os quartos de desempenho ).
Segundo o desempenho, observam -se diferenças estatisticamente significativas entre o quarto de desempenho inferior e os demais.
No quarto superior, a clareza e objetividade de todos ou da maioria dos enunciados das questões foram percebidas por 84,5 \%. 4.3.2 Componente de Conhecimento Específico Em relação aos enunciados das questões do Componente de Conhecimento Específico da prova, para 85,6 \% dos estudantes avaliados da Área de Engenharia Florestal, a clareza e a objetividade ( Questão 5 ) estavam presentes em todas ou na maioria das questões ( Gráfico 4.9, Gráfico 4.10, e no Anexo II, a Tabela II.5 ).
A maioria dos estudantes de todas as Grandes Regiões brasileiras consideraram claros e objetivos todos ou a maioria dos enunciados das questões do Componente de Conhecimento Específico da prova, percentual sempre maior ou igual a 84,0 \%.
Não são observadas diferenças estatisticamente significativas entre as Grandes Regiões.
A proporção de estudantes que consideraram os enunciados das questões claros e objetivos apresenta uma tendência crescente em relação ao aumento de desempenho : mais elevada no quarto superior ( 91,0 \% ), se comparada ao quarto inferior de desempenho ( 79,1 \% ).
Observa -se que as diferenças entre o quarto inferior e os dois quartos superiores de desempenho são estatisticamente significativas. 4.4 SUFICIÊNCIA DAS INFORMAÇÕES/INSTRUÇÕES FORNECIDAS Ao avaliarem as informações/instruções fornecidas para a resolução das questões ( Questão 6 ), 89,2 \% dos respondentes da Área de Engenharia Florestal de todo o Brasil afirmaram que eram até excessivas ou suficientes em todas ou na maioria das questões ( Gráfico 4.11, Gráfico 4.12, e, no Anexo II, a Tabela II.6 ).
Quanto à distribuição de respondentes pelas Grandes Regiões observa -se que a proporção de estudantes que consideraram as informações/instruções fornecidas até excessivas ou suficientes em todas ou na maioria das questões foi sempre superior ou igual a 86,7 \%, percentual da região Nordeste, chegando a 91,1 \% na região Centro-Oeste.
As diferenças entre as proporções das regiões não são estatisticamente significativas.
Levando -se em conta o desempenho dos participantes, notam -se diferenças estatisticamente significativas entre as opiniões de estudantes do quarto superior de desempenho em relação ao primeiro e ao terceiro quartos, como mostrado no Gráfico 4.12.
O percentual de participantes que avaliaram as informações/instruções como até excessivas ou suficientes em todas ou na maioria das questões foi mais elevado no quarto superior ( 93,6 \% ), percentual superior à média nacional ( 89,2 \% ).
Já no quarto inferior, as informações/instruções avaliadas pelos respondentes como até excessivas, em todas ou na maioria das questões, atingiram o percentual de 85,2 \%. 4.5 DIFICULDADE ENCONTRADA AO RESPONDER À PROVA Indagados sobre as dificuldades com as quais se depararam ao responder à prova ( Questão 7 ), 20,2 \% dos estudantes apontaram o Desconhecimento do conteúdo.
Para 46,1 \%, a Forma diferente de abordagem do conteúdo foi indicada como dificuldade.
Já a Falta de motivação para fazer a prova foi a dificuldade apontada por 12,9 \% dos respondentes.
Considerando -se todo o Brasil, 17,3 \% dos respondentes afirmaram que não tiveram qualquer tipo de dificuldade para responder à prova ( Tabela II.7 no Anexo II ).
Os Gráficos 4.13 e 4.14 apresentam os percentuais de estudantes que apontaram o Desconhecimento do conteúdo como dificuldade ao responder à prova.
Na análise por Grandes Regiões, o percentual de inscritos e presentes que apontaram o Desconhecimento do conteúdo como dificuldade ao responder à prova não superou 25,3 \%.
Os percentuais variaram de 17,4 \% na região Sudeste a 25,3 \% na região Nordeste, não havendo diferença significativa estatisticamente entre nenhuma das regiões.
A Forma diferente de abordagem do conteúdo foi a escolha modal dos estudantes, com percentuais que variaram de 40,6 \% ( região Centro-Oeste ) a 50,1 \% ( região Sul ).
O percentual de alunos que citaram a Falta de motivação como dificuldade variou de 7,8 \% ( região Nordeste ) a 16,0 \% ( região Centro-Oeste ).
Os que declararam não ter qualquer tipo de dificuldade para responder à prova variaram de 14,7 \% na região Norte a 20,1 \% na região Sudeste.
Em relação aos quartos de desempenho, o Desconhecimento do conteúdo foi a opção escolhida por 18,2 \% dos estudantes do quarto superior e por 21,5 \% dos estudantes do quarto inferior, não havendo diferenças estatisticamente significativas dentre os quartos de desempenho.
A alternativa modal para os alunos, quando agregados pelos quartos de desempenho, foi a de que a dificuldade encontrada foi causada pela Forma diferente de abordagem do conteúdo : 44,3 \% no quarto inferior e 46,5 \% do quarto superior escolheram essa alternativa. 4.6 CONTEÚDOS DAS QUESTÕES OBJETIVAS DA PROVA Ao analisar os conteúdos das questões objetivas da prova ( Questão 8 ), um percentual muito pequeno dos estudantes avaliados, apenas 3,1 \%, afirmou que não estudou ainda a maioria desses conteúdos ( Gráficos 4.15, Gráfico 4.16, e a Tabelas II.8 no Anexo II ).
A maioria ( 75,7 \% ) afirmou ter estudado e aprendido muitos ou todos os conteúdos.
Na análise por Grande Região, a proporção de respondentes que escolheram a opção não estudou ainda a maioria desses conteúdos foi pequena.
Observa -se que nas regiões Norte ( 4,7 \% ) e Nordeste ( 6,1 \% ), apesar de pequenas, as proporções foram maiores do que a média nacional ( 3,1 \% ).
Tanto a região Sul quanto a região Sudeste apresentam diferenças estatisticamente significativas em relação às regiões Norte e Nordeste.
Em todas as regiões, a maioria dos presentes afirmou ter estudado e aprendido muitos ou todos os conteúdos, com proporções variando entre 71,0 \% na região Nordeste e 85,0 \% na região Sudeste.
Considerando -se separadamente as opiniões de estudantes dos quatro quartos de desempenho, observa -se que, no quarto inferior, 6,2 \% ofereceram como resposta que não estudou ainda a maioria desses conteúdos, sendo 1,0 \% os do quarto superior com a mesma resposta, observando -se diferença estatisticamente significativa entre os alunos que optaram por esse motivo de dificuldade no quarto inferior de desempenho em relação aos dois quartos superiores.
Levando -se em conta o quarto superior, 90,6 \% dos alunos afirmaram ter estudado e aprendido muitos ou todos os conteúdos.
No outro extremo, no primeiro quarto, 60,2 \% optaram pelas mesmas categorias.
As proporções são crescentes com o desempenho. 4.7 TEMPO GASTO PARA CONCLUIR A PROVA Ao responderem sobre o tempo de conclusão da prova ( Questão 9 ), mais de três quartos dos estudantes ( 79,1 \% ) afirmaram ter gasto Entre duas e quatro horas ( Gráfico 4.17, Gráfico 4.18 e, no Anexo II, a Tabela II.9 ).
Considerando -se as cinco Grandes Regiões brasileiras, os percentuais dos que utilizaram Entre duas e quatro horas para finalizar a prova nas regiões Sul ( 74,5 \% ) e Centro-Oeste ( 72,2 \% ) foram inferiores ao percentual nacional.
Nas demais Grandes Regiões, o percentual de alunos que dispensaram Entre duas e quatro horas para concluir a prova ficou igual ou acima de 80,0 \%, como mostra o Gráfico 4.17.
Há diferenças estatisticamente significativas entre a proporção do Sudeste quando comparada à proporção registrada nas regiões Sul e Centro-Oeste.
Uma vez analisadas as alternativas escolhidas pelos estudantes que se situam nos diferentes quartos de desempenho, verifica -se que uma proporção maior de participantes no quarto superior ( 85,4 \% ) declarou ter gasto Entre duas e quatro horas para concluir a prova em relação aos quartos inferiores.
As diferenças do primeiro quarto de desempenho ( 71,7 \% ) em relação ao segundo e ao quarto são estatisticamente significativas.

% matched lemmas: 
\end{textsample}
